
Estimates of public opinion are critically important for assessing representation, accountability, and democratic legitimacy. But inferences based on comparing constituency-level opinion to policy outcomes require accurate, fine-grained measures of public opinion — measures that are increasingly difficult to obtain given declining response rates and growing nonignorable nonresponse. As a result, it is often unclear whether discrepancies between public opinion and policy reflect failures of representation or failures of polling.

We propose a method that limits the impact of survey errors on small-area opinion estimates by leveraging auxiliary data with known ground truth.
Building on the logit shift calibration framework common in MRP, we extend the approach to multiple correlated outcomes.
The key insight is that discrepancies between model-based and ground-truth estimates for auxiliary variables (e.g., election results) are informative about likely survey errors for correlated outcomes for which no  ground truth exists (e.g., policy attitudes). 
We propose a principled calibration method building on this intuition. 
The calibration is driven by geographic random effect correlations estimated from a joint model of all outcomes.
The method produces larger adjustments for more highly correlated variables. 
We also show how this approach can be extended to hyperlocal geographies lacking poststratification data, such as precincts, by calibrating to ground-truth margins at the lower level (Appendix \ref{apx:precinct}).

Using pre-election polling from the 2022 Michigan midterm, we demonstrate that calibrating county-level MRP estimates to gubernatorial results reduces error by as much as two-thirds. 
We find similar gains when extending the approach to precinct-level results.


A key benefit of our approach is that it can correct for multiple sources of survey error --- nonignorable nonresponse, model misspecification, and errors in poststratification data --- without requiring researchers to decompose the contribution of each source. 
As demonstrated by simulations in Appendix \ref{apx:simulations}, our calibration procedure is most beneficial when the calibration variable is strongly correlated with the outcome of interest and when nonignorable nonresponse is present. 
Even with small, highly nonrepresentative samples, calibration substantially reduces error for correlated outcomes.

The method relies on the covariance equivalence assumption, which states that the geographic correlation structure estimated from the survey reflects the population-level correlations. 
In Appendix \ref{apx:covariance-assumption}, we demonstrate the plausibility of this assumption in our validation setting.
Of course, this assumption may not always hold, particularly if survey respondents' attitudes are more correlated than those of nonrespondents. 
However, the alternative --- standard MRP without calibration --- implicitly assumes that observed errors in auxiliary variables are completely uninformative about errors in other outcomes.
Our approach is likely to improve estimates even when the assumption is not perfectly met.

By providing researchers with a principled method for incorporating known population quantities into MRP estimates, we help ensure that substantive conclusions based on small-area opinion measures are robust to survey error. 
The method complements existing MRP tools. 
Finally, our method can be implemented using our accompanying \texttt{R} package, \texttt{calibratedMRP}.
